% GLOBAL GLOSSARY



% ------------------------------------------------
% GLOSSARIES

\newglossaryentry{markupsprache}{
    name={Markupsprache},
    description={reichert Texte mit zusätzlichen Informationen an}}
\newglossaryentry{latex}{
    name={LaTeX},
    description={ist eine \Gls{markupsprache}, die häufig in der Wissenschaft und Technik verwendet wird}}


\newglossaryentry{VonRestorffEffekt}{
    name={Von-Restorff-Effekt},
    description={Mnemonisches Phänomen, bei dem sich von ihrer Umgebung abweichende Reize besser einprägen}
}

\newglossaryentry{FittschesGesetz}{
    name={Fitts'sche Gesetz},
    description={Mathematische Formel zur Vorhersage der Zeit, die benötigt wird, um ein Ziel zu erreichen, basierend auf Distanz und Zielgröße}
}

\newglossaryentry{HickschesGesetz}{
    name={Hick'sche Gesetz},
    description={Psychophysikalisches Modell zur Vorhersage von Reaktionszeiten bei Entscheidungen mit multiplen Wahloptionen}
}

\newglossaryentry{progressiveDisclosure}{
    name={Progressive-Disclosure-Muster},
    description={Gestaltungsprinzip, das Informationen schrittweise offenbart, um kognitive Überlastung zu vermeiden}
}

\newglossaryentry{hedonicpragmaticmodel}{
    name={Hedonic-Pragmatic-Model},
    description={Designkonzept, das sowohl die sinnliche Freude (hedonisch) als auch die funktionale Nutzbarkeit (pragmatisch) eines Produkts berücksichtigt}
}

\newglossaryentry{eyecatcher}{
    name = {Eyecatcher},
    description= {Visuelles Element, das gezielt eingesetzt wird, um die Aufmerksamkeit der Betrachter zu erregen und deren Fokus auf eine bestimmte Information oder ein Detail zu lenken}
}

\newglossaryentry{usability}{
    name={Usability},
    description={Bezeichnet die Gebrauchstauglichkeit eines Produkts in Bezug auf Effektivität, Effizienz und Zufriedenheit der Nutzung. Im UX-Design ist eine hohe Usability ein zentrales Ziel}
}

\newglossaryentry{wcag}{
    name={WCAG-Kontrastverhältnisse},
    description={Helligkeitsunterschied zwischen Text/Vordergrund und Hintergrund gemäß den \Gls{WCAG-abbr}}
}
\newglossaryentry{SM}{
    name={Systemmodell},
    description={Technische Implementierungslogiken aus Designerperspektive}
}

\newglossaryentry{UT}{
    name={Usability Testing},
    description={Systematische Evaluation der Benutzerfreundlichkeit eines Produkts durch reale Nutzende}
}

\newglossaryentry{prototyping}{
    name={Prototyping},
    description={Iterativer Prozess zur Erstellung vereinfachter Modelle, um Designkonzepte frühzeitig zu testen und zu optimieren}
}

\newglossaryentry{Farbpsychologie}{
    name={Farbpsychologie},
    description={Wissenschaft, die untersucht, wie Farben menschliches Verhalten und Emotionen beeinflussen}
}

\newglossaryentry{Microinteractions}{
    name = {Microinteractions},
    description = {Kleine, oft unsichtbare Interaktionen in digitalen Produkten, die Nutzerfeedback geben}
}

\newglossaryentry{Farbharmonien}{
    name={Farbharmonien},
    description={Ästhetisch kohärente Farbkombinationen zur Schaffung visueller Konsistenz}
}

\newglossaryentry{Interfaces}{
    name={Interfaces},
    description={
            Schnittstellen zwischen Mensch und Maschine, die Interaktionen ermöglichen}
}
    
\newglossaryentry{TaskCompletionRaten}{
    name={Task Completion Raten},
    description={
            Kennzahl in Usability Tests, die angibt, wie viele Nutzende eine gestellte Aufgabe erfolgreich und vollständig abschließen konnten. Dient als Indikator für die Verständlichkeit und Bedienbarkeit eines Interfaces}
}

% ------------------------------------------------
% ACRONYMS


\newacronym[shortplural=Abken, longplural=Abkürzungen]{abk}{Abk}{Abkürzung}

\newacronym{etal}{et al.}{und andere}
\newacronym{iAa}{i. A. a.}{in Anlehnung an}
\newacronym{vgl}{vgl.}{vergleiche}
\newacronym{na}{N/A}{not available / nicht angegeben}
\newacronym{ua}{u. a.}{unter anderem}

\newacronym[shortplural=RQs, longplural=Forschungsfragen]{rq}{RQ}{Forschungsfrage}
\newacronym{ux}{UX}{User Experience}
\newacronym{ui}{UI}{User Interface}
\newacronym{gui}{GUI}{Graphical User Interface}
\newacronym{uxd}{UX-Design}{User Experience Design}
\newacronym{uiux}{UI/UX}{User Interface/User Experience}
\newacronym{uxresearch}{UX-Research}{User Experience Research}
\newacronym{clt}{CLT}{Cognitve Load Theory}
\newacronym{cta}{CTA}{Call-to-Action}
\newacronym{WCAG-abbr}{WCAG}{Web Content Accessibility Guidelines}